% SOURCE_STATEMENT: "Is $\\sum \\frac{(5n)!}{(n!)^5}z^n$ related to (non-holomorphic) forms? Is it related to automorphic objects?\nDo other solutions to the Picard-Fuchs equation have modular properties?"
% Preparation and final self-audit: CONTRIBUTING.md prompts C and D; source/audit artifacts retained in the local physics_statements_2026_10_11 batch.
% !TeX program = lualatex
\documentclass[11pt]{article}
\usepackage[margin=25mm]{geometry}
\usepackage{fontspec}
\setmainfont{Latin Modern Roman}
\usepackage{amsmath,amssymb,mathtools}
\usepackage{unicode-math}
\setmathfont{Latin Modern Math}
\usepackage{xurl,hyperref}
\hypersetup{colorlinks=true,urlcolor=blue}
\setcounter{secnumdepth}{0}
\setlength{\parindent}{0pt}
\setlength{\parskip}{5pt}
\setlength{\emergencystretch}{3em}
\newcommand{\sref}[2]{\hyperlink{source:#1:#2}{[S#2]}}
\newcommand{\cataloguescope}[1]{\paragraph{Recorded target.}#1}
\newcommand{\R}{\mathbb R}
\newcommand{\C}{\mathbb C}
\newcommand{\N}{\mathbb N}
\newcommand{\Z}{\mathbb Z}
\newcommand{\Q}{\mathbb Q}
\newcommand{\F}{\mathbb F}
\newcommand{\eps}{\varepsilon}
\begin{document}
% UnsolvedMath ID 20002513; source catalogue code AIM-PHYSICS-0021
\setcounter{section}{20002512}
\section{Mirror-quintic periods, modular-type structures, and a conifold obstruction}\label{20002513}
{\small UnsolvedMath ID 20002513 · Mathematical Physics}
\subsection{Definitions and mathematical statement}

\paragraph{Shared notation.}$\mathbb N=\{1,2,\ldots\}$; the standard number fields have their usual meanings. All additional conventions are stated locally.


\paragraph{Definitions and notation.}The distinguished series is $\varpi_0(z)=\sum_{n\ge0}(5n)!/(n!)^5\,z^n$. With $\theta=z\,d/dz$, it satisfies the mirror-quintic Picard--Fuchs operator $\theta^4-5z(5\theta+1)(5\theta+2)(5\theta+3)(5\theta+4)$. Other independent solutions describe additional periods, possibly with logarithmic behavior. Modular or automorphic behavior requires a group, domain and transformation law.

\paragraph{Catalogue wording retained for traceability.}
Is $\sum \frac{(5n)!}{(n!)^5}z^n$ related to (non-holomorphic) forms? Is it related to automorphic objects?
Do other solutions to the Picard-Fuchs equation have modular properties?
\paragraph{Scope and source qualification.}The Picard–Fuchs equation identifies the series without claiming a classical modular uniformization. The source asks about nonholomorphic forms and other automorphic objects as well as the distinguished period; all are retained. The catalogue title/status was paired with a machine-generated research attempt; that attempt is not peer-reviewed evidence of settlement. The statement covers the original source questions rather than only the title’s subcase.
\subsection{Short English statement}
Determine modular or automorphic structures for the mirror-quintic period series and the other Picard–Fuchs solutions.
\subsection{Sources}
\begin{itemize}
\item[\hypertarget{source:20002513:1}{S1}] ULAM.AI and the UnsolvedMath Contributors, UnsolvedMath, record 20002513 (AIM-PHYSICS-0021). Catalogue curation under CC BY 4.0; definitions and English exposition added locally.
\url{https://huggingface.co/datasets/ulamai/UnsolvedMath}.
\item[\hypertarget{source:20002513:2}{S2}] American Institute of Mathematics, Gromov-Witten invariants and number theory, source item 4.4.
\url{http://aimpl.org/gromwitnumthry/4/}.
\end{itemize}
\label{end:20002513}
\end{document}
